\subsectionfit{Location}
\label{config:location}
\index{location}
\index{Config!location}

\begin{lstlisting}[language=Go,style=kaolstplain]
location( "name" "latitude" "longitude" ["altitude"] )
\end{lstlisting}

The \Statement{location} statement defines the location on the Earth's surface
where a weather station is located.
It's used to enable the system to perform various calculations like when the Sun
rises or sets as well as other phenomena specific to your location.

\ParameterDescription{
    \item[name] the unique name\sidenote{As the system can manage multiple weather stations, name must be unique.
    This uniqueness can be just locally, but if you join multiple systems together then this must be unique across all of them.
    } of the station.
    \item[latitude] latitude of the station
    \item[longitude] longitude of the station
    \item[altitude] altitude\sidenote{
        Altitude is optional, however advisable to set it as it is used in certain calculations like
        calculating the Pressure at mean sea level and in forecasts.
    } of the station in meters above mean sea level.
}

Latitude and Longitude are defined in one of the following formats:

\begin{table}[hb]
    \raggedright
    \begin{tabular}{ l|l }
        \hrule
            12.345      & Degrees in decimal \\
            12:34.56    & Degrees \& decimal Minutes \\
            12:34:56.78 & Degrees, Minutes \& decimal Seconds \\
        \hrule
    \end{tabular}
\end{table}

For declaring which hemisphere the station is in, you can prefix the Degrees field with one of the prefixes
in the following table.
If no prefix is included then the location is in the North/East respectively.

\begin{table}[hb]
    \begin{tabular}{ c c c }
        \toprule
        Hemisphere & Latitude & Longitude \\
        \midrule
        Northern    & + N   & \\
        Southern    & - S   & \\
        \midrule
        Eastern     &       & + E \\
        Western     &       & - W \\
        \bottomrule
    \end{tabular}
\end{table}

